\documentclass[a4paper,12pt]{article}

\usepackage[T2A]{fontenc}
\usepackage[utf8]{inputenc}
\usepackage[russian]{babel}

\usepackage{cmap}
\usepackage{amsmath,amssymb}
\usepackage{geometry}
\usepackage{graphicx}
\usepackage[most]{tcolorbox}
\usepackage{fancyhdr}
\usepackage[hidelinks]{hyperref}

\geometry{margin=2cm}
\setlength{\headheight}{17pt}

\pagestyle{fancy}
\fancyhf{}
\fancyhead[L]{ЛАФНП. ЭКЗАМЕН. 2-й СЕМЕСТР}
\fancyhead[R]{Артём Ким ИУ7-22Б}
\fancyfoot[C]{\thepage}

\newtcolorbox{definitionbox}{
colback=white,
colframe=green!60!black,
boxrule=0pt,
leftrule=3pt,
rightrule=0pt,
toprule=0pt,
bottomrule=0pt,
left=6pt,
right=0pt,
top=3pt,
bottom=3pt
}

\newtcolorbox{theorembox}{
colback=white,
colframe=red!70!pink,
boxrule=0pt,
leftrule=3pt,
rightrule=0pt,
toprule=0pt,
bottomrule=0pt,
left=6pt,
right=0pt,
top=3pt,
bottom=3pt
}

\newtcolorbox{notebox}{
colback=white,
colframe=blue!65!cyan,
boxrule=0pt,
leftrule=3pt,
rightrule=0pt,
toprule=0pt,
bottomrule=0pt,
left=6pt,
right=0pt,
top=3pt,
bottom=3pt
}

\begin{document}

\pagenumbering{gobble}

\begin{titlepage}
    \begin{center}
        \large Московский государственный технический университет\\
        имени Н.\,Э. Баумана\\[1cm]

        \textbf{2-й семестр}\\[5cm]

        {\LARGE \textbf{КОНСПЕКТ}}\\[0.5cm]
        {\Large лекций}\\[2cm]

        {\large по дисциплине:}\\[0.3cm]
        {\Large \textbf{Линейная алгебра и функции нескольких переменных}}\\[3cm]
    \end{center}

    \vfill

    \begin{center}
        2026
    \end{center}
\end{titlepage}

\clearpage
\pagenumbering{arabic}

\tableofcontents
\newpage

\section{Функции многих переменных. Предел и непрерывность}

\subsection*{Функции многих переменных}

Рассмотрим отображение

\[
f\colon D \subset \mathbb{R}^n \to \mathbb{R}
\]

\begin{definitionbox}
Функцией многих переменных называется отображение, которое каждой точке
$x = (x_1, x_2, \dots, x_n) \in D \subset \mathbb{R}^n$
ставит в соответствие единственное число
$f(x_1, x_2, \dots, x_n) \in \mathbb{R}$.
\end{definitionbox}

Множество $D$ называется областью определения функции.

Геометрически функция двух переменных задаёт поверхность:

\[
z = f(x,y)
\]

\subsection*{Окрестность точки}

\begin{definitionbox}
Окрестностью точки $a \in \mathbb{R}^n$ называется множество точек, расстояние от которых до точки $a$ меньше некоторого числа $\varepsilon > 0$:
\end{definitionbox}

\[
U_\varepsilon(a) =
\{x \in \mathbb{R}^n \mid |x-a| < \varepsilon\}
\]

\subsection*{Предел функции многих переменных}

\begin{definitionbox}
Число $A$ называется пределом функции $f(x)$ при $x \to a$, если
\end{definitionbox}

\[
\forall \varepsilon > 0 \;
\exists \delta > 0\colon
0 < |x-a| < \delta
\Rightarrow
|f(x)-A| < \varepsilon
\]

Обозначение:

\[
\lim_{x \to a} f(x) = A
\]

\subsection*{Непрерывность функции}

\begin{definitionbox}
Функция $f$ называется непрерывной в точке $a$, если
\end{definitionbox}

\[
\lim_{x \to a} f(x) = f(a)
\]

То есть

\[
\forall \varepsilon > 0\;
\exists \delta > 0\colon
|x-a| < \delta
\Rightarrow
|f(x)-f(a)| < \varepsilon
\]

\subsection*{Последовательностный критерий предела}

\begin{theorembox}
\[
\lim_{x \to a} f(x) = A
\]

тогда и только тогда, когда для любой последовательности

\[
x_k \to a,
\qquad x_k \ne a
\]

выполняется

\[
f(x_k) \to A
\]
\end{theorembox}

\subsection*{Пример функции}

\[
f(x,y) = \sqrt{x^2 + y^2}
\]

Область определения:

\[
D = \mathbb{R}^2
\]

\subsection*{Пример предела}

\[
\lim\limits_{(x,y)\to(0,0)} \sqrt{x^2 + y^2} = 0
\]

\end{document}